\documentclass[aspectratio=169, 12pt]{beamer}

% ── Pakete ──────────────────────────────────────────────────
\usepackage[ngerman]{babel}
\usepackage[utf8]{inputenc}
\usepackage[T1]{fontenc}
\usepackage{amsmath, amssymb, mathtools}
\usepackage{tikz}
\usetikzlibrary{arrows.meta, calc, positioning, decorations.markings,
                 shapes.geometric, backgrounds, fit, matrix, angles, quotes}
\usepackage{graphicx}
\usepackage{booktabs}
\usepackage{fontawesome5}
\usepackage{hyperref}
\usepackage[table]{xcolor}
\usepackage{multimedia}
\usepackage{animate}

% ── Theme ───────────────────────────────────────────────────
\usetheme{metropolis}
\setbeamercolor{frametitle}{bg=black!85, fg=white}
\setbeamercolor{progress bar}{fg=orange}
\metroset{progressbar=frametitle, block=fill}

% ── Farben ──────────────────────────────────────────────────
\definecolor{akzent}{HTML}{E8590C}
\definecolor{hintergrund}{HTML}{FAF9F6}
\definecolor{dunkel}{HTML}{1A1A2E}
\definecolor{mittel}{HTML}{16213E}
\definecolor{hell}{HTML}{0F3460}
\definecolor{highlight}{HTML}{E94560}
\setbeamercolor{background canvas}{bg=hintergrund}

% ── Hilfsbefehle ────────────────────────────────────────────
\newcommand{\placeholder}[2][0.5\textwidth]{%
  \begin{center}
    \fcolorbox{akzent}{white!95!black}{%
      \parbox{#1}{\centering\vspace{1.0cm}%
        {\color{akzent}\faImage\quad\textit{#2}}%
        \vspace{1.0cm}}}
  \end{center}
}

\newcommand{\stich}[1]{\item[\color{akzent}\faAngleRight] #1}

\newcommand{\formelbox}[1]{%
  \begin{center}
    \colorbox{dunkel!10}{%
      \parbox{0.88\textwidth}{%
        \centering
        $\displaystyle #1$%
      }%
    }
  \end{center}
}

\newcommand{\abschnitt}[2]{%
  {
  \setbeamercolor{background canvas}{bg=dunkel}
  \begin{frame}[plain]
    \vfill
    \begin{center}
      {\Huge\color{white}\textbf{#1}}\\[1em]
      {\large\color{white!70}#2}
    \end{center}
    \vfill
  \end{frame}
  }
}

\newcommand{\quelle}[1]{{\vfill\hfill\scriptsize\color{gray}#1}}

% ── Titelseite ──────────────────────────────────────────────
\title{\textbf{Wie KI rechnen lernt}}
\subtitle{Was passiert in einem neuronalen Netz, wenn es Aufgaben l\"ost?}
\author{}
\date{}

% ============================================================
\begin{document}
% ============================================================

\begin{frame}[plain]
  \maketitle
\end{frame}

% ════════════════════════════════════════════════════════════
\abschnitt{Die Frage}{Kann ein neuronales Netz einen Algorithmus \emph{erfinden}?}
% ════════════════════════════════════════════════════════════

\begin{frame}{Worum es geht}
  \begin{center}
    \colorbox{dunkel!10}{%
      \parbox{0.85\textwidth}{\centering\large
        Wie lernt eine KI, etwas N\"utzliches zu tun?\\[4pt]
        Und \textbf{wie genau} f\"uhrt sie es dann aus?
      }%
    }
  \end{center}

  \bigskip
  \begin{itemize}
    \stich{KI-Modelle sehen Beispiele und erkennen Muster}
    \stich{Aber was \emph{genau} passiert dabei im Inneren?}
    \stich{Wir schauen uns ein konkretes, einfaches Beispiel an,\\
           bei dem Forscher den gelernten Algorithmus \textbf{vollst\"andig aufgedeckt} haben}
    \stich{Das hilft uns zu verstehen, wie KI generell funktioniert}
  \end{itemize}

  \quelle{Nanda, Chan, Lieberum, Smith, Steinhardt. ``Progress Measures for Grokking via Mechanistic Interpretability.'' ICLR 2023.}
\end{frame}

\begin{frame}{Das Experiment}
  \begin{center}
    {\Large\texttt{\color{dunkel} 37 + 58 = \color{highlight}\textbf{?} \quad (mod 113)}}
  \end{center}

  \bigskip
  \begin{itemize}
    \stich{Ein winziges neuronales Netz bekommt Rechenaufgaben}
    \stich{Es sieht nur Zahlenpaare und die richtige Antwort}
    \stich{Niemand erkl\"art ihm, \emph{wie} man rechnet}
    \stich{Frage: Lernt es \textbf{auswendig} -- oder entdeckt es einen \textbf{Trick}?}
  \end{itemize}

  \bigskip
  \begin{center}
    \colorbox{highlight!10}{%
      \parbox{0.75\textwidth}{\centering\small
        Spoiler: Es erfindet einen eleganten Algorithmus,\\
        den ihm niemand beigebracht hat.
      }%
    }
  \end{center}
\end{frame}

% ════════════════════════════════════════════════════════════
\abschnitt{Modulo-Rechnung}{Was bedeutet ``mod''?}
% ════════════════════════════════════════════════════════════

\begin{frame}{Was ist Modulo? -- Konkrete Beispiele}
  \begin{center}
    {\large Modulo = \textbf{Rest bei der Division}}
  \end{center}

  \bigskip
  \begin{columns}[T]
    \column{0.48\textwidth}
    \textbf{Einfache Beispiele (mod 12):}
    \begin{itemize}
      \stich{$5 + 3 = 8$ \quad {\color{gray}(bleibt 8)}}
      \stich{$9 + 5 = 14 \to \mathbf{2}$ \quad {\color{gray}(14 $-$ 12)}}
      \stich{$11 + 4 = 15 \to \mathbf{3}$ \quad {\color{gray}(15 $-$ 12)}}
    \end{itemize}

    \bigskip
    {\small\color{gray} Wie eine Uhr: nach 12 geht es\\wieder bei 0 los!}

    \column{0.48\textwidth}
    \textbf{Im Experiment (mod 113):}
    \begin{itemize}
      \stich{$37 + 58 = 95$ \quad {\color{gray}(bleibt 95)}}
      \stich{$80 + 50 = 130 \to \mathbf{17}$ \quad {\color{gray}(130 $-$ 113)}}
      \stich{$100 + 100 = 200 \to \mathbf{87}$ \quad {\color{gray}(200 $-$ 113)}}
    \end{itemize}

    \bigskip
    {\small\color{gray} Gleiche Idee, nur mit 113\\statt 12.}
  \end{columns}
\end{frame}

\begin{frame}{Warum Modulo auf einem Kreis lebt}
  \begin{columns}[T]
    \column{0.45\textwidth}
    \begin{center}
      \begin{tikzpicture}[scale=1.6]
        \draw[gray!50, thick] (0,0) circle (1);
        
        % Uhrziffern
        \foreach \h/\angle in {0/90, 3/0, 6/270, 9/180} {
          \node[font=\small, gray] at (\angle:1.25) {\h};
        }
        
        % Markierungen
        \foreach \h in {0,...,11} {
          \pgfmathsetmacro{\angle}{90 - \h*30}
          \draw[gray!60] (\angle:0.92) -- (\angle:1.0);
        }
        
        % Beispiel: 9 + 5 = 2 (mod 12)
        \pgfmathsetmacro{\startangle}{90 - 9*30}
        \pgfmathsetmacro{\endangle}{90 - 14*30}
        
        \fill[blue!70] (\startangle:1) circle (2.5pt);
        \node[blue!70, font=\scriptsize\bfseries] at (\startangle:1.4) {9};
        
        \draw[akzent, thick, -{Stealth}] (\startangle:0.75) arc (\startangle:\endangle:0.75);
        \node[akzent, font=\scriptsize] at (90 - 11.5*30:0.5) {$+5$};
        
        \pgfmathsetmacro{\resultangle}{90 - 2*30}
        \fill[green!60!black] (\resultangle:1) circle (3pt);
        \node[green!60!black, font=\scriptsize\bfseries] at (\resultangle:1.4) {2};
      \end{tikzpicture}
    \end{center}

    \column{0.52\textwidth}
    \begin{itemize}
      \stich{Eine Uhr ist ein Kreis mit 12 Positionen}
      \stich{$9 + 5$: Starte bei 9, gehe 5 Schritte weiter}
      \stich{Nach 12 kommt wieder 0 $\to$ Ergebnis: \textbf{2}}
      \stich{Addition auf dem Kreis = \textbf{Rotation}!}
    \end{itemize}

    \bigskip
    \colorbox{akzent!10}{%
      \parbox{0.9\linewidth}{\centering\small
        \textbf{Kernidee:}\\
        Modulo-Addition ist nichts anderes\\
        als \textbf{Drehen auf einem Kreis}.
      }%
    }
  \end{columns}
\end{frame}

\begin{frame}{Vom Uhren-Kreis zum Einheitskreis}
  \begin{center}
    \begin{tikzpicture}[scale=1.4]
      % Einheitskreis
      \draw[gray!50, thick] (0,0) circle (1);
      \draw[-{Stealth}, gray!30] (-1.4,0) -- (1.4,0) node[right, font=\tiny]{x};
      \draw[-{Stealth}, gray!30] (0,-1.4) -- (0,1.4) node[above, font=\tiny]{y};
      
      % Zahl a = 37, Winkel = 37 * (360/113) ≈ 117.9°
      \pgfmathsetmacro{\angA}{37*360/113}
      \fill[blue!70] (\angA:1) circle (2.5pt);
      \node[blue!70, font=\small\bfseries] at (\angA:1.3) {$37$};
      \draw[blue!70, dashed] (0,0) -- (\angA:1);
      
      % Zahl b = 58, Winkel = 58 * (360/113) ≈ 184.8°
      \pgfmathsetmacro{\angB}{58*360/113}
      
      % Rotation von a um b
      \pgfmathsetmacro{\angSum}{95*360/113}
      \fill[green!60!black] (\angSum:1) circle (3pt);
      \node[green!60!black, font=\small\bfseries] at (\angSum:1.35) {$95$};
      \draw[green!60!black, thick] (0,0) -- (\angSum:1);
      
      % Rotationspfeil (im Uhrzeigersinn von a nach a+b)
      \draw[akzent, thick, -{Stealth}] (\angA:0.7) arc (\angA:\angSum:0.7);
      \node[akzent, font=\scriptsize] at ({(\angA+\angSum)/2}:0.45) {$+58$};
      
      % Annotation
      \node[font=\scriptsize, gray, text width=4cm, align=center] at (3.2, 0) {
        Jede Zahl $t$ sitzt bei\\
        Winkel $\frac{t}{113} \times 360°$\\[6pt]
        $37 + 58 = 95$\\
        $(95 < 113$, also bleibt es 95)
      };
    \end{tikzpicture}
  \end{center}

  \bigskip
  \begin{center}
    \colorbox{dunkel!10}{%
      \parbox{0.8\textwidth}{\centering
        Alle 113 Zahlen sitzen gleichm\"a\ss ig auf einem Kreis.\\[3pt]
        \textbf{Addition = Weiterdrehen.} Nach 113 Schritten: zur\"uck zum Start.
      }%
    }
  \end{center}
\end{frame}

% ════════════════════════════════════════════════════════════
\abschnitt{Grokking}{Pl\"otzliches Verstehen -- nach langem Auswendiglernen}
% ════════════════════════════════════════════════════════════

\begin{frame}{Was ist Grokking?}
  \begin{columns}[T]
    \column{0.55\textwidth}
    \begin{itemize}
      \stich{Das Netz lernt die Trainingsbeispiele\\
             \textbf{sofort auswendig} (100\% richtig)}
      \stich{Bei \emph{neuen} Aufgaben: \textbf{lange 0\% richtig}}
      \stich{Nach tausenden weiteren Runden:\\
             pl\"otzlich \textbf{alles richtig} -- auch Neues!}
      \stich{Es hat einen \textbf{Algorithmus entdeckt}}
    \end{itemize}

    \bigskip
    {\small\color{gray} Benannt nach Robert Heinlein:\\
    ``to grok'' = tiefes, intuitives Verstehen}

    \column{0.4\textwidth}
    \placeholder[0.9\linewidth]{Bild: Train/Test Accuracy Kurve\\
    (Figure 2 aus Nanda et al.\ 2023)}
  \end{columns}

  \quelle{Power et al.\ 2022 (Erstentdeckung); Nanda et al.\ 2023 (Erkl\"arung)}
\end{frame}

\begin{frame}{Grokking -- die drei Phasen}
  \begin{center}
    \begin{tikzpicture}[scale=0.9]
      \draw[-{Stealth}, thick] (0,0) -- (13,0) node[right]{\small Trainingsrunden};
      
      \fill[akzent!15] (0,0.3) rectangle (3,2.5);
      \fill[hell!15] (3,0.3) rectangle (8,2.5);
      \fill[highlight!15] (8,0.3) rectangle (12,2.5);
      
      \node[font=\small\bfseries, text=akzent] at (1.5,2.0) {Phase 1};
      \node[font=\scriptsize, text width=2.5cm, align=center] at (1.5,1.2) {Auswendiglernen\\(wie Vokabeln pauken)};
      
      \node[font=\small\bfseries, text=hell] at (5.5,2.0) {Phase 2};
      \node[font=\scriptsize, text width=3cm, align=center] at (5.5,1.2) {Algorithmus bildet sich\\(im Hintergrund!)};
      
      \node[font=\small\bfseries, text=highlight] at (10,2.0) {Phase 3};
      \node[font=\scriptsize, text width=2.5cm, align=center] at (10,1.2) {Aufr\"aumen\\(Auswendiggelerntes\\wird gel\"oscht)};
      
      \draw[dashed, gray] (3,0) -- (3,2.5);
      \draw[dashed, gray] (8,0) -- (8,2.5);
    \end{tikzpicture}
  \end{center}

  \bigskip
  \begin{itemize}
    \stich{\textbf{\"Uberraschung:} Der Algorithmus ist \emph{l\"angst fertig}, bevor das Netz ihn benutzt!}
    \stich{Erst wenn das Auswendiglernen ``teurer'' wird als der Algorithmus, schaltet es um}
  \end{itemize}
\end{frame}

% ════════════════════════════════════════════════════════════
\abschnitt{Der Trick des Netzes}{Wie es Addition in Drehung verwandelt}
% ════════════════════════════════════════════════════════════

\begin{frame}{Die Grundidee: Zahlen als Punkte auf Kreisen}
  \begin{center}
    \colorbox{dunkel!10}{%
      \parbox{0.85\textwidth}{\centering\large
        Das Netz legt jede Zahl auf einen \textbf{Kreis}\\[4pt]
        und verwandelt Addition in \textbf{Drehung}.
      }%
    }
  \end{center}

  \bigskip
  \begin{center}
    \begin{tikzpicture}[scale=1.4]
      \draw[gray!50, thick] (0,0) circle (1);
      \draw[-{Stealth}, gray!30] (-1.3,0) -- (1.3,0);
      \draw[-{Stealth}, gray!30] (0,-1.3) -- (0,1.3);
      
      % Zahl a
      \pgfmathsetmacro{\angA}{37*360*14/113}
      \pgfmathsetmacro{\angAmod}{mod(\angA,360)}
      \fill[blue!70] (\angAmod:1) circle (2.5pt);
      \node[blue!70, font=\small\bfseries] at (\angAmod:1.3) {$a$};
      \draw[blue!70, dashed] (0,0) -- (\angAmod:1);
      
      % Ergebnis a+b
      \pgfmathsetmacro{\angS}{95*360*14/113}
      \pgfmathsetmacro{\angSmod}{mod(\angS,360)}
      \fill[green!60!black] (\angSmod:1) circle (3pt);
      \node[green!60!black, font=\small\bfseries] at (\angSmod:1.35) {$a+b$};
      \draw[green!60!black, thick] (0,0) -- (\angSmod:1);
      
      % Rotationspfeil
      \draw[akzent, thick, -{Stealth}] (\angAmod:0.65) arc (\angAmod:\angSmod:0.65);
      \node[akzent, font=\scriptsize] at ({(\angAmod+\angSmod)/2}:0.4) {$+b$};
    \end{tikzpicture}
  \end{center}

  \begin{center}
    \small Starte bei $a$, drehe um den Winkel von $b$ weiter $\to$ lande bei $a+b$ (mod 113).
  \end{center}
\end{frame}

\begin{frame}{Warum Sinus und Cosinus?}
  \begin{columns}[T]
    \column{0.45\textwidth}
    \begin{center}
      \begin{tikzpicture}[scale=1.8]
        \draw[gray!50, thick] (0,0) circle (1);
        \draw[-{Stealth}, gray!30] (-1.2,0) -- (1.2,0) node[right, font=\tiny]{cos};
        \draw[-{Stealth}, gray!30] (0,-1.2) -- (0,1.2) node[above, font=\tiny]{sin};
        
        % Punkt bei Winkel 50°
        \pgfmathsetmacro{\ang}{50}
        \fill[akzent] (\ang:1) circle (2.5pt);
        \draw[akzent, thick] (0,0) -- (\ang:1);
        
        % Projektion cos
        \draw[blue!70, thick, dashed] (\ang:1) -- ({cos(\ang)},0);
        \fill[blue!70] ({cos(\ang)},0) circle (2pt);
        \node[blue!70, font=\scriptsize, below] at ({cos(\ang)}, -0.1) {$\cos(\alpha)$};
        
        % Projektion sin
        \draw[red!70, thick, dashed] (\ang:1) -- (0,{sin(\ang)});
        \fill[red!70] (0,{sin(\ang)}) circle (2pt);
        \node[red!70, font=\scriptsize, left] at (-0.1, {sin(\ang)}) {$\sin(\alpha)$};
        
        % Winkel
        \draw[gray, thick] (0.3,0) arc (0:\ang:0.3);
        \node[gray, font=\scriptsize] at (25:0.45) {$\alpha$};
      \end{tikzpicture}
    \end{center}

    \column{0.52\textwidth}
    \begin{itemize}
      \stich{Ein Punkt auf dem Kreis hat zwei Koordinaten: \textbf{cos} und \textbf{sin}}
      \stich{Zusammen beschreiben sie die \textbf{Position} auf dem Kreis}
      \stich{Das Netz speichert jede Zahl $t$ als:\\[4pt]
             {\color{blue!70}$\cos(\alpha)$} und {\color{red!70}$\sin(\alpha)$}\\[4pt]
             wobei $\alpha = \frac{2\pi k \cdot t}{113}$}
      \stich{Das ist die ``Adresse'' der Zahl auf dem Kreis}
    \end{itemize}
  \end{columns}
\end{frame}

\begin{frame}<1>[label=embedding1]{Schritt 1: Embedding -- Zahlen auf den Kreis legen}
  \begin{center}
    {\large Frequenz $k = 14$: Jede Zahl bekommt ihren Platz}
  \end{center}

  \bigskip
  \begin{center}
    \begin{tikzpicture}[scale=1.8]
      \draw[gray!40, thick] (0,0) circle (1);
      \draw[-{Stealth}, gray!30] (-1.2,0) -- (1.2,0) node[right, font=\tiny]{cos};
      \draw[-{Stealth}, gray!30] (0,-1.2) -- (0,1.2) node[above, font=\tiny]{sin};
      
      % Nur wenige Zahlen zeigen, gut verteilt
      \pgfmathsetmacro{\angZero}{mod(0*360*14/113,360)}
      \fill[black] (\angZero:1) circle (2pt);
      \node[font=\scriptsize, black] at (\angZero:1.25) {$0$};
      
      \pgfmathsetmacro{\angTen}{mod(10*360*14/113,360)}
      \fill[blue!70] (\angTen:1) circle (2pt);
      \node[font=\scriptsize, blue!70] at (\angTen:1.25) {$10$};
      
      \pgfmathsetmacro{\angThirty}{mod(30*360*14/113,360)}
      \fill[red!70] (\angThirty:1) circle (2pt);
      \node[font=\scriptsize, red!70] at (\angThirty:1.25) {$30$};
      
      \pgfmathsetmacro{\angFifty}{mod(50*360*14/113,360)}
      \fill[green!60!black] (\angFifty:1) circle (2pt);
      \node[font=\scriptsize, green!60!black] at (\angFifty:1.25) {$50$};
      
      \pgfmathsetmacro{\angEighty}{mod(80*360*14/113,360)}
      \fill[violet] (\angEighty:1) circle (2pt);
      \node[font=\scriptsize, violet] at (\angEighty:1.25) {$80$};
      
      \pgfmathsetmacro{\angHundred}{mod(100*360*14/113,360)}
      \fill[orange] (\angHundred:1) circle (2pt);
      \node[font=\scriptsize, orange] at (\angHundred:1.25) {$100$};
    \end{tikzpicture}
  \end{center}

  \begin{center}
    \small Zahl $t$ $\to$ Winkel $\frac{2\pi \cdot 14 \cdot t}{113}$ $\to$ Punkt auf dem Kreis.\\
    Verschiedene Zahlen landen an verschiedenen Stellen.
  \end{center}
\end{frame}

\begin{frame}<2>[label=embedding2]{Schritt 1: Embedding -- Alle 113 Zahlen}
  \begin{center}
    {\large Alle Zahlen von 0 bis 112 auf dem Kreis (Frequenz $k=14$)}
  \end{center}

  \bigskip
  \begin{center}
    \begin{tikzpicture}[scale=1.8]
      \draw[gray!40, thick] (0,0) circle (1);
      \draw[-{Stealth}, gray!30] (-1.2,0) -- (1.2,0);
      \draw[-{Stealth}, gray!30] (0,-1.2) -- (0,1.2);
      
      % Alle 113 Zahlen
      \foreach \t in {0,1,...,112} {
        \pgfmathsetmacro{\angle}{mod(\t*360*14/113,360)}
        \fill[akzent!70] (\angle:1) circle (0.8pt);
      }
      
      \node[font=\small, below] at (0,-1.4) {113 Punkte, gleichm\"a\ss ig verteilt};
    \end{tikzpicture}
  \end{center}

  \begin{center}
    \small Da 14 und 113 keinen gemeinsamen Teiler haben,\\
    landen alle 113 Zahlen an \textbf{verschiedenen} Stellen -- der Kreis wird gleichm\"a\ss ig gef\"ullt.
  \end{center}
\end{frame}

\begin{frame}{Mehrere Kreise gleichzeitig}
  \begin{center}
    {\large Das Netz benutzt \textbf{mehrere Frequenzen} -- also mehrere Kreise}
  \end{center}

  \bigskip
  \begin{center}
    \begin{tikzpicture}[scale=1.0]
      % Kreis 1
      \begin{scope}[xshift=0cm]
        \draw[gray!40, thick] (0,0) circle (1);
        % Zeige Zahl 10 auf diesem Kreis
        \pgfmathsetmacro{\angA}{mod(10*360*14/113,360)}
        \fill[blue!70] (\angA:1) circle (2.5pt);
        \node[font=\tiny, blue!70] at (\angA:1.3) {10};
        % Zeige Zahl 20
        \pgfmathsetmacro{\angB}{mod(20*360*14/113,360)}
        \fill[red!70] (\angB:1) circle (2.5pt);
        \node[font=\tiny, red!70] at (\angB:1.3) {20};
        \node[font=\scriptsize, below] at (0,-1.4) {Kreis 1 ($k=14$)};
        \node[font=\tiny, gray, above] at (0,1.3) {langsame Drehung};
      \end{scope}
      
      % Kreis 2
      \begin{scope}[xshift=3.5cm]
        \draw[gray!40, thick] (0,0) circle (1);
        \pgfmathsetmacro{\angA}{mod(10*360*35/113,360)}
        \fill[blue!70] (\angA:1) circle (2.5pt);
        \node[font=\tiny, blue!70] at (\angA:1.3) {10};
        \pgfmathsetmacro{\angB}{mod(20*360*35/113,360)}
        \fill[red!70] (\angB:1) circle (2.5pt);
        \node[font=\tiny, red!70] at (\angB:1.3) {20};
        \node[font=\scriptsize, below] at (0,-1.4) {Kreis 2 ($k=35$)};
        \node[font=\tiny, gray, above] at (0,1.3) {mittlere Drehung};
      \end{scope}
      
      % Kreis 3
      \begin{scope}[xshift=7cm]
        \draw[gray!40, thick] (0,0) circle (1);
        \pgfmathsetmacro{\angA}{mod(10*360*52/113,360)}
        \fill[blue!70] (\angA:1) circle (2.5pt);
        \node[font=\tiny, blue!70] at (\angA:1.3) {10};
        \pgfmathsetmacro{\angB}{mod(20*360*52/113,360)}
        \fill[red!70] (\angB:1) circle (2.5pt);
        \node[font=\tiny, red!70] at (\angB:1.3) {20};
        \node[font=\scriptsize, below] at (0,-1.4) {Kreis 3 ($k=52$)};
        \node[font=\tiny, gray, above] at (0,1.3) {schnelle Drehung};
      \end{scope}
    \end{tikzpicture}
  \end{center}

  \bigskip
  \begin{itemize}
    \stich{Dieselbe Zahl sitzt auf \textbf{jedem Kreis an einer anderen Stelle}}
    \stich{H\"ohere Frequenz $k$ = gr\"o\ss erer Winkelschritt pro Zahl}
    \stich{Das Netz benutzt \textbf{5 Kreise gleichzeitig} -- warum, sehen wir gleich}
  \end{itemize}
\end{frame}

\begin{frame}{Warum mehrere Kreise? -- Das Problem mit einem einzigen}
  \begin{columns}[T]
    \column{0.48\textwidth}
    \begin{center}
      {\small\textbf{Nur 1 Frequenz ($k=14$):}}\\[0.5em]
      \begin{tikzpicture}[scale=0.7]
        \draw[-{Stealth}, gray!30] (-0.5,0) -- (6,0) node[right, font=\tiny]{$c$};
        \draw[-{Stealth}, gray!30] (0,-1.5) -- (0,1.5) node[above, font=\tiny]{Wert};
        
        \draw[blue!70, thick, domain=0:5.5, samples=80]
          plot (\x, {cos(360*14*\x/113*10)});
        
        % Markiere Maximum
        \fill[green!60!black] (0, 1) circle (3pt);
        \node[font=\tiny, green!60!black, above right] at (0,1) {richtig!};
        
        % Markiere Nebenpeak
        \pgfmathsetmacro{\nebenpeak}{113/14*0.5}
        \fill[red!70] (4.04, {cos(360*14*4.04/113*10)}) circle (3pt);
        \node[font=\tiny, red!70, above] at (4.04, 0.8) {fast gleich hoch!};
        
        \node[font=\scriptsize, below, text width=4cm, align=center] at (2.75,-1.8) 
          {Problem: Andere Werte von $c$\\sind \textbf{fast genauso gro\ss}};
      \end{tikzpicture}
    \end{center}

    \column{0.48\textwidth}
    \begin{center}
      {\small\textbf{5 Frequenzen zusammen:}}\\[0.5em]
      \begin{tikzpicture}[scale=0.7]
        \draw[-{Stealth}, gray!30] (-0.5,0) -- (6,0) node[right, font=\tiny]{$c$};
        \draw[-{Stealth}, gray!30] (0,-2) -- (0,3.5) node[above, font=\tiny]{Wert};
        
        \draw[akzent, line width=2pt, domain=0:5.5, samples=100]
          plot (\x, {cos(360*14*\x/113*10) + cos(360*35*\x/113*10) 
                    + cos(360*41*\x/113*10) + cos(360*42*\x/113*10)
                    + cos(360*52*\x/113*10)});
        
        % Markiere Maximum
        \fill[green!60!black] (0, 3) circle (4pt);
        \node[font=\tiny, green!60!black, above right] at (0,3) {$c = a+b$: alle $=1$!};
        
        \node[font=\scriptsize, below, text width=4cm, align=center] at (2.75,-2.3) 
          {L\"osung: Nebenpeaks der verschiedenen\\Frequenzen \textbf{liegen woanders}\\$\to$ l\"oschen sich gegenseitig aus};
      \end{tikzpicture}
    \end{center}
  \end{columns}

  \bigskip
  \begin{center}
    \colorbox{akzent!10}{%
      \parbox{0.8\textwidth}{\centering\small
        \textbf{Konstruktive Interferenz:} Nur beim richtigen $c$ zeigen alle 5 Wellen nach oben.\\
        \"Uberall sonst zeigen sie in verschiedene Richtungen und heben sich auf.
      }%
    }
  \end{center}
\end{frame}

% ════════════════════════════════════════════════════════════
\abschnitt{Schritt f\"ur Schritt}{Was passiert im Netz?}
% ════════════════════════════════════════════════════════════

\begin{frame}{Schritt 2: Attention -- Informationen zusammenbringen}
  \begin{center}
    \begin{tikzpicture}[
        box/.style={draw, thick, rounded corners, minimum width=2cm,
                     minimum height=0.8cm, font=\small},
        arr/.style={-{Stealth}, thick}
      ]
      \node[box, fill=blue!15] (a) at (0,0) {Position 1: $a$};
      \node[box, fill=red!15] (b) at (5,0) {Position 2: $b$};
      \node[box, fill=gray!15] (eq) at (10,0) {Position 3: $=$};
      
      \draw[arr, blue!60, thick] (a.north) to[out=50, in=130] 
        node[above, font=\scriptsize]{kopiert $\cos(\omega_k a)$, $\sin(\omega_k a)$} (eq.north);
      \draw[arr, red!60, thick] (b.north) to[out=50, in=130] 
        node[below, font=\scriptsize, yshift=3pt]{kopiert $\cos(\omega_k b)$, $\sin(\omega_k b)$} (eq.north east);
    \end{tikzpicture}
  \end{center}

  \bigskip
  \begin{itemize}
    \stich{Das Netz liest das Ergebnis an Position ``$=$'' ab}
    \stich{Attention \textbf{kopiert} die Kreis-Koordinaten von $a$ und $b$ dorthin}
    \stich{Zus\"atzlich: Attention-Gewichte sind selbst periodisch\\
           $\to$ erzeugen \textbf{Produkte} wie $\cos(\omega_k a) \cdot \cos(\omega_k b)$}
    \stich{Genau diese Produkte braucht man f\"ur den n\"achsten Schritt!}
  \end{itemize}
\end{frame}

\begin{frame}{Schritt 3: MLP -- Die Additionsformel berechnen}
  \begin{center}
    {\large Wie wird aus $\cos(\omega a)$ und $\cos(\omega b)$ ein $\cos(\omega(a+b))$?}
  \end{center}

  \bigskip
  \begin{center}
    \colorbox{dunkel!10}{%
      \parbox{0.85\textwidth}{\centering
        \textbf{Additionstheorem} (Schulwissen):\\[6pt]
        $\cos(\alpha + \beta) = \cos\alpha \cdot \cos\beta \;-\; \sin\alpha \cdot \sin\beta$
      }%
    }
  \end{center}

  \bigskip
  \begin{itemize}
    \stich{Das Netz muss also \textbf{Produkte} berechnen: $\cos(\omega a) \times \cos(\omega b)$}
    \stich{Problem: Neuronen k\"onnen nur \textbf{addieren} und \textbf{ReLU} anwenden}
    \stich{Aber: Viele ReLU-Neuronen zusammen k\"onnen Produkte \textbf{ann\"ahern}!}
  \end{itemize}
\end{frame}

\begin{frame}{Wie ReLU Produkte ann\"ahern kann}
  \begin{center}
    \begin{tikzpicture}[scale=0.9]
      % ReLU Funktion
      \draw[-{Stealth}, gray!50] (-2.5,0) -- (3,0) node[right, font=\small]{$x$};
      \draw[-{Stealth}, gray!50] (0,-0.5) -- (0,3) node[above, font=\small]{ReLU($x$)};
      
      \draw[blue!70, line width=2pt] (-2.5,0) -- (0,0);
      \draw[akzent, line width=2pt] (0,0) -- (2.5,2.5);
      
      \fill[black] (0,0) circle (2pt);
      \node[font=\scriptsize, below right] at (0.2,-0.3) {Knick bei 0};
      
      % Annotation rechts
      \node[font=\small, text width=5cm, align=left] at (7, 1.5) {
        \textbf{Der Trick:}\\[4pt]
        $\text{ReLU}(a+b)^2 - \text{ReLU}(a-b)^2$\\[2pt]
        $\approx$ Vielfaches von $a \times b$\\[8pt]
        512 Neuronen mit verschiedenen\\
        Gewichten n\"ahern gemeinsam\\
        die ben\"otigten Produkte an.\\[8pt]
        {\color{gray}\scriptsize (>90\% der Varianz erkl\"art)}
      };
    \end{tikzpicture}
  \end{center}

  \bigskip
  \begin{center}
    \colorbox{akzent!10}{%
      \parbox{0.75\textwidth}{\centering\small
        \textbf{Ergebnis:} Die MLP-Schicht gibt $\cos(\omega_k(a+b))$ und $\sin(\omega_k(a+b))$ aus --\\
        die \textbf{Kreis-Koordinaten der Summe}!
      }%
    }
  \end{center}
\end{frame}

\begin{frame}{Schritt 4: Unembedding -- Das Ergebnis ablesen}
  \begin{center}
    {\large F\"ur jede m\"ogliche Antwort $c$ (von 0 bis 112):}
  \end{center}

  \bigskip
  \begin{center}
    \begin{tikzpicture}[scale=1.2]
      \draw[gray!40, thick] (0,0) circle (1);
      \draw[-{Stealth}, gray!30] (-1.3,0) -- (1.3,0);
      \draw[-{Stealth}, gray!30] (0,-1.3) -- (0,1.3);
      
      % a+b Punkt
      \pgfmathsetmacro{\angSum}{95*360*14/113}
      \pgfmathsetmacro{\angSumMod}{mod(\angSum,360)}
      \fill[akzent] (\angSumMod:1) circle (3pt);
      \node[font=\scriptsize, akzent] at (\angSumMod:1.35) {$a+b$};
      \draw[akzent, thick] (0,0) -- (\angSumMod:1);
      
      % c Punkt (richtig)
      \fill[green!60!black] (\angSumMod:1) circle (3pt);
      \node[font=\scriptsize, green!60!black, below left] at (\angSumMod:1.4) {$c^* = a+b$};
      
      % c Punkt (falsch)
      \pgfmathsetmacro{\angWrong}{50*360*14/113}
      \pgfmathsetmacro{\angWrongMod}{mod(\angWrong,360)}
      \fill[gray!60] (\angWrongMod:1) circle (2pt);
      \node[font=\tiny, gray] at (\angWrongMod:1.3) {anderes $c$};
      
      % Annotation
      \node[font=\small, text width=5cm, align=left] at (3.5, 0) {
        Vergleiche den Winkel von $a+b$\\
        mit dem Winkel von jedem $c$:\\[6pt]
        \textbf{Gleicher Winkel} $\to$ $\cos(0) = 1$\\
        {\color{green!60!black}$\to$ hoher Score!}\\[4pt]
        \textbf{Anderer Winkel} $\to$ $\cos(\neq 0) < 1$\\
        {\color{gray}$\to$ niedriger Score}
      };
    \end{tikzpicture}
  \end{center}

  \bigskip
  \begin{center}
    \small Das Netz berechnet: Wie \"ahnlich ist die Position von $c$ zur Position von $a+b$?\\
    $\to$ Je \"ahnlicher, desto h\"oher der Score (``Logit'') f\"ur diese Antwort.
  \end{center}
\end{frame}

% ════════════════════════════════════════════════════════════
\abschnitt{Konstruktive Interferenz}{Wie 5 Wellen das Ergebnis finden}
% ════════════════════════════════════════════════════════════

\begin{frame}{Schritt 5: Die Wellen \"uberlagern sich}
  \begin{center}
    \begin{tikzpicture}[scale=0.75,
      declare function={
        fA(\x) = cos(360*14*\x/113);
        fB(\x) = cos(360*35*\x/113);
        fC(\x) = cos(360*41*\x/113);
        fD(\x) = cos(360*42*\x/113);
        fE(\x) = cos(360*52*\x/113);
      }]
      
      % Achsen
      \draw[-{Stealth}, gray!30] (-0.5,0) -- (10,0) node[right, font=\small]{$c$};
      \draw[-{Stealth}, gray!30] (0,-1.5) -- (0,1.5);
      
      % Einzelne Wellen (d\"unn, halbtransparent)
      \draw[blue!40, thin, domain=0:9.5, samples=100]
        plot (\x, {0.25*fA(\x)});
      \draw[red!40, thin, domain=0:9.5, samples=100]
        plot (\x, {0.25*fB(\x)});
      \draw[green!40!black, thin, opacity=0.5, domain=0:9.5, samples=100]
        plot (\x, {0.25*fC(\x)});
      \draw[violet!40, thin, domain=0:9.5, samples=100]
        plot (\x, {0.25*fD(\x)});
      \draw[orange!40, thin, domain=0:9.5, samples=100]
        plot (\x, {0.25*fE(\x)});
      
      % Summe (dick, gestrichelt)
      \draw[akzent, line width=2.5pt, dashed, domain=0:9.5, samples=120]
        plot (\x, {0.25*(fA(\x) + fB(\x) + fC(\x) + fD(\x) + fE(\x))});
      
      % Peak markieren
      \draw[akzent, thick, -{Stealth}] (0, 1.6) -- (0, 1.35);
      \node[akzent, font=\small\bfseries, above] at (0, 1.6) {$c = a+b$};
      \node[akzent, font=\tiny] at (0, 1.9) {Alle 5 Wellen $= 1$ hier!};
      
      % Legende
      \node[font=\tiny, blue!60] at (8.5, 1.4) {$k=14$};
      \node[font=\tiny, red!60] at (8.5, 1.1) {$k=35$};
      \node[font=\tiny, green!50!black] at (8.5, 0.8) {$k=41$};
      \node[font=\tiny, violet] at (8.5, 0.5) {$k=42$};
      \node[font=\tiny, orange] at (8.5, 0.2) {$k=52$};
      \node[font=\tiny\bfseries, akzent] at (8.5, -0.2) {Summe};
    \end{tikzpicture}
  \end{center}

  \bigskip
  \begin{itemize}
    \stich{Jede d\"unne Linie = eine Frequenz (Cosinus-Welle)}
    \stich{Die \textbf{gestrichelte Summe} hat einen klaren Peak bei $c = a+b$}
    \stich{\"Uberall sonst: die Wellen zeigen in verschiedene Richtungen $\to$ heben sich auf}
  \end{itemize}
\end{frame}

\begin{frame}{Zusammenfassung: Der ganze Algorithmus}
  \begin{center}
    \begin{tikzpicture}[
        box/.style={draw, thick, rounded corners, minimum width=2.5cm,
                     minimum height=1.2cm, font=\scriptsize, align=center},
        arr/.style={-{Stealth}, thick, akzent},
        node distance=0.6cm
      ]
      \node[box, fill=hell!15] (emb) {
        \textbf{1. Embedding}\\[2pt]
        Zahl $t$ $\to$\\
        Punkt auf 5 Kreisen
      };
      \node[box, fill=dunkel!15, right=of emb] (attn) {
        \textbf{2. Attention}\\[2pt]
        Kopiert Kreis-\\
        Koordinaten zusammen
      };
      \node[box, fill=akzent!15, right=of attn] (mlp) {
        \textbf{3. MLP}\\[2pt]
        Berechnet\\
        $\cos(\omega_k(a+b))$
      };
      \node[box, fill=highlight!15, right=of mlp] (unemb) {
        \textbf{4. Ablesen}\\[2pt]
        Vergleicht mit\\
        jedem $c$ $\to$ Score
      };
      
      \draw[arr] (emb) -- (attn);
      \draw[arr] (attn) -- (mlp);
      \draw[arr] (mlp) -- (unemb);
    \end{tikzpicture}
  \end{center}

  \bigskip
  \begin{center}
    \colorbox{dunkel!10}{%
      \parbox{0.85\textwidth}{\centering
        \textbf{In Worten:}\\[4pt]
        Lege die Zahlen auf Kreise. Drehe den Kreis um $b$ weiter.\\
        Schaue, welche Antwort $c$ am n\"achsten am Ergebnis liegt.\\[4pt]
        5 Kreise gleichzeitig $\to$ eindeutiges Ergebnis.
      }%
    }
  \end{center}

  \bigskip
  \begin{center}
    \small\color{gray} Das Netz hat diesen Algorithmus \textbf{selbst erfunden} -- 
    niemand hat ihm Trigonometrie beigebracht!
  \end{center}
\end{frame}

% ════════════════════════════════════════════════════════════
\abschnitt{Evidenz}{Woher wissen wir, dass das stimmt?}
% ════════════════════════════════════════════════════════════

\begin{frame}{Beweis 1: Die Gewichte sind periodisch}
  \begin{columns}[T]
    \column{0.48\textwidth}
    \placeholder[0.9\linewidth]{Bild: Fourier-Komponenten von $W_E$\\
    (Figure 3 links aus Nanda et al.)\\[4pt]
    \footnotesize Nur 5 von 56 Frequenzen\\haben nennenswerte St\"arke}

    \bigskip
    \begin{center}
      \small\textbf{Embedding-Matrix:}\\
      Fast alle Energie in 5 Frequenzen\\
      $k \in \{14, 35, 41, 42, 52\}$
    \end{center}

    \column{0.48\textwidth}
    \placeholder[0.9\linewidth]{Bild: Neuron-Aktivierungen\\
    (Figure 4 Mitte aus Nanda et al.)\\[4pt]
    \footnotesize Klare periodische Muster}

    \bigskip
    \begin{center}
      \small\textbf{Neuron-Aktivierungen:}\\
      Jedes Neuron reagiert periodisch\\
      auf die Eingaben
    \end{center}
  \end{columns}

  \quelle{Nanda et al.\ 2023, Figures 3 \& 4}
\end{frame}

\begin{frame}{Beweis 2: Ablation -- Entfernen best\"atigt den Algorithmus}
  \begin{center}
    \begin{tikzpicture}[scale=0.85]
      % Drei Szenarien
      \node[font=\small\bfseries] at (0, 3.5) {Original};
      \node[font=\small\bfseries] at (4, 3.5) {Nur 5 Frequenzen};
      \node[font=\small\bfseries] at (8.5, 3.5) {5 Frequenzen entfernt};
      
      % Balken
      \fill[hell!60] (-0.5,0) rectangle (0.5, 1.5);
      \node[font=\tiny] at (0, -0.3) {Loss: $2.4 \cdot 10^{-7}$};
      \node[font=\scriptsize, hell] at (0, 1.8) {gut};
      
      \fill[green!50] (3.5,0) rectangle (4.5, 1.0);
      \node[font=\tiny] at (4, -0.3) {Loss: $7.2 \cdot 10^{-8}$};
      \node[font=\scriptsize, green!60!black] at (4, 1.3) {\textbf{besser!}};
      
      \fill[highlight!60] (8,0) rectangle (9, 3.2);
      \node[font=\tiny] at (8.5, -0.3) {Loss: $5.27$};
      \node[font=\scriptsize, highlight] at (8.5, 3.5) {\textbf{kaputt}};
      
      % Achse
      \draw[gray!50] (-1,0) -- (10,0);
    \end{tikzpicture}
  \end{center}

  \bigskip
  \begin{itemize}
    \stich{95\% der Frequenzen entfernen, nur die 5 wichtigen behalten: \textbf{wird sogar besser!}}
    \stich{Die 5 wichtigen Frequenzen entfernen: \textbf{schlechter als Raten}}
    \stich{$\Rightarrow$ Der Algorithmus steckt \textbf{genau} in diesen 5 Frequenzen}
  \end{itemize}

  \quelle{Nanda et al.\ 2023, Figure 6 \& Section 4.4}
\end{frame}

% ════════════════════════════════════════════════════════════
\abschnitt{Analogie}{Was bedeutet das f\"ur KI allgemein?}
% ════════════════════════════════════════════════════════════

\begin{frame}{Analogie: Geografie lernen}
  \begin{columns}[T]
    \column{0.55\textwidth}
    \begin{itemize}
      \stich{Gurnee \& Tegmark (2023):\\
             LLMs lernen \textbf{Weltkarten}!}
      \stich{Nie explizit trainiert --\\
             emergiert als beste Kompression}
      \stich{\"Ahnlich wie hier: das Netz entdeckt\\
             \textbf{Kreise} als beste Darstellung\\
             f\"ur modulare Arithmetik}
      \stich{In beiden F\"allen: das System lernt\\
             \textbf{geometrische Strukturen}\\
             eigenst\"andig}
    \end{itemize}

    \column{0.4\textwidth}
    \begin{center}
      \placeholder[0.85\linewidth]{Bild: Weltkarte aus LLM-Embeddings\\
      (Gurnee \& Tegmark 2023)\\[4pt]
      \footnotesize 2D-Projektion zeigt\\erkannte Geografie}
    \end{center}
  \end{columns}

  \quelle{Gurnee, Tegmark. ``Language Models Represent Space and Time.'' arXiv:2310.02207 (2023)}
\end{frame}

\begin{frame}{Kreise hier -- Karten dort}
  \begin{columns}[T]
    \column{0.45\textwidth}
    \begin{center}
      \textbf{Modulare Addition}\\[0.5em]
      \begin{tikzpicture}[scale=1.0]
        \draw[gray!40, thick] (0,0) circle (1);
        \foreach \t in {0,8,...,112} {
          \pgfmathsetmacro{\angle}{mod(360*14*\t/113,360)}
          \fill[akzent] (\angle:1) circle (1.2pt);
        }
        \node[font=\scriptsize] at (0,-1.5) {Kreise = beste Darstellung};
      \end{tikzpicture}
    \end{center}

    \column{0.45\textwidth}
    \begin{center}
      \textbf{Sprache \"uber Orte}\\[0.5em]
      \begin{tikzpicture}[scale=1.0]
        \draw[gray!40, thick] (0,0) ellipse (1.5 and 1);
        \fill[hell!60] (-0.7, 0.4) circle (3pt) node[font=\tiny, above]{Paris};
        \fill[hell!60] (0.3, 0.3) circle (3pt) node[font=\tiny, above]{Berlin};
        \fill[hell!60] (0.9, -0.1) circle (3pt) node[font=\tiny, right]{Wien};
        \fill[akzent] (-0.9, -0.4) circle (3pt) node[font=\tiny, below]{Madrid};
        \fill[akzent] (0.6, 0.6) circle (3pt) node[font=\tiny, above]{Warschau};
        \node[font=\scriptsize] at (0,-1.5) {Karte = beste Darstellung};
      \end{tikzpicture}
    \end{center}
  \end{columns}

  \bigskip
  \begin{center}
    \colorbox{dunkel!10}{%
      \parbox{0.85\textwidth}{\centering
        \textbf{Gemeinsames Prinzip:}\\[4pt]
        KI-Systeme entdecken eigenst\"andig die geometrische Struktur,\\
        die am besten zu ihren Daten passt -- ohne dass man es ihnen sagt.
      }%
    }
  \end{center}
\end{frame}

% ════════════════════════════════════════════════════════════
\abschnitt{Einordnung}{Vom kleinen Baustein zum gro\ss en Modell}
% ════════════════════════════════════════════════════════════

\begin{frame}{Ein winziger Baustein in einem riesigen System}
  \begin{center}
    \begin{tikzpicture}[scale=0.9]
      % Großes Modell als Rahmen
      \draw[thick, dunkel, rounded corners, fill=dunkel!5] (-5,-2.5) rectangle (5,2.5);
      \node[font=\small\bfseries, dunkel] at (0,2.8) {Gro\ss es Sprachmodell};
      
      % Circuits als kleine Kreise, h\"ubsch angeordnet
      \foreach \i/\label/\col in {
        1/Grammatik/blue!50, 2/Fakten/green!50, 3/Logik/violet!50,
        4/Rechnen/akzent, 5/Analogien/orange!50, 6/Kontext/teal!50} {
        \pgfmathsetmacro{\xpos}{-3.5 + mod(\i-1,3)*3.5}
        \pgfmathsetmacro{\ypos}{1.0 - floor((\i-1)/3)*2.2}
        \node[draw=\col, thick, fill=\col!20, rounded corners, 
              minimum width=2.2cm, minimum height=              minimum height=1cm, font=\scriptsize, align=center}
          ] at (\xpos, \ypos) {\label};
      }
      
      % Highlight "Rechnen"
      \node[draw=akzent, line width=2pt, fill=akzent!30, rounded corners,
            minimum width=2.2cm, minimum height=1cm, font=\scriptsize\bfseries,
            align=center] at (0, 1.0) {Rechnen};
      
      % Pfeil zum Rechnen-Knoten
      \draw[akzent, thick, -{Stealth}] (0, -0.2) -- (0, -1.8)
        node[below, font=\scriptsize, text width=4cm, align=center, text=akzent]
        {\textbf{Dieser Vortrag:}\\ein einzelner, winziger Baustein};
    \end{tikzpicture}
  \end{center}

  \bigskip
  \begin{itemize}
    \stich{Gro\ss e Sprachmodelle bestehen aus \textbf{hunderten solcher Bausteine}}
    \stich{Jeder Baustein l\"ost eine andere Aufgabe}
    \stich{Wir haben \textbf{einen einzigen} davon vollst\"andig verstanden}
  \end{itemize}
\end{frame}

\begin{frame}{Was wir gelernt haben}
  \begin{enumerate}
    \item[\color{akzent}\textbf{1.}] KI-Systeme k\"onnen \textbf{eigenst\"andig Algorithmen erfinden}
    \item[\color{akzent}\textbf{2.}] Sie nutzen daf\"ur \textbf{geometrische Strukturen} (hier: Kreise)
    \item[\color{akzent}\textbf{3.}] \textbf{Grokking:} Erst auswendig lernen, dann pl\"otzlich verstehen
    \item[\color{akzent}\textbf{4.}] \textbf{Konstruktive Interferenz:} Mehrere Wellen finden gemeinsam die Antwort
    \item[\color{akzent}\textbf{5.}] Wir k\"onnen in die Black Box \textbf{hineinschauen} und verstehen, was passiert
  \end{enumerate}

  \bigskip
  \begin{center}
    \colorbox{dunkel!10}{%
      \parbox{0.85\textwidth}{\centering
        \textbf{Die gro\ss e Erkenntnis:}\\[4pt]
        Neuronale Netze sind keine magischen Black Boxes.\\
        Sie lernen konkrete, nachvollziehbare Algorithmen --\\
        man muss nur genau genug hinschauen.
      }%
    }
  \end{center}
\end{frame}

% ════════════════════════════════════════════════════════════
% Schlussfolie
% ════════════════════════════════════════════════════════════

\begin{frame}[plain]
  {
  \setbeamercolor{background canvas}{bg=dunkel}
  \begin{tikzpicture}[remember picture, overlay]
    \node at (current page.center) {%
      \begin{minipage}{0.8\textwidth}
        \centering
        {\Huge\color{white}\textbf{Danke!}}\\[1.5em]
        {\large\color{white!70}Fragen?}\\[2em]
        {\small\color{white!50}
          Paper: Nanda et al.\ ``Progress Measures for Grokking\\
          via Mechanistic Interpretability'' (ICLR 2023)\\[0.5em]
          Code \& interaktive Figuren: \texttt{neelnanda.io/grokking-paper}
        }
      \end{minipage}
    };
  \end{tikzpicture}
  }
\end{frame}

% ============================================================
\end{document}
% ============================================================
